%lezione 20 del 6 maggio 2026 pag 45 appunti
%-----------------------------------------------------------------------------------------------------

Esiste anche una condizione a basse energie in cui $|c(k)|^2<<1$: se $kr\approx0$,
$\sin(kr)\approx kr$ e l'integrale diventa
$$|c(k)| = \frac{2m}{\hbar^2k}\int_{0}^{\infty}dr|kr||V(r)| = \frac{2m}{\hbar^2}\int_{0}^{\infty}dr
r|V(r)|\leq\frac{2m}{\hbar^2}V_0\int_{0}^{A}dr r = \frac{mV_0A^2}{\hbar^2}<<1
$$
che impone la condizione
$$V_0<<\frac{\hbar^2}{mA^2}$$
\textbf{Oss.} Se il potenziale è attrattivo, questa condizione implica che esso non può
formare stati legati: utilizzando il fatto che $pA\sim\hbar$ si ha che
$$E<0\leftrightarrow \frac{p^2}{2m}+V \approx \frac{\hbar^2}{2mA^2}-V_0<0 \leftrightarrow
V_0 > \frac{\hbar^2}{2mA^2}$$
la formazione di stati legati è dunque incompatibile con l'approssimazione di Born per basse energie.
\subsection{Esempio: il potenziale di Yukawa}
Il potenziale di Yukawa è un potenziale centrale a range finito, definito come
$$V(r)=A\frac{e^{-\mu r}}{r}$$
con $A\in\mathbb{R},\phantom{x}\mu\in\mathbb{R}^+$ costanti: il termine esponenziale negativo
garantisce che il potenziale svanisca molto rapidamente lontano da $r=0$.

Per alte energie si può dunque applicare l'approssimazione di Born, calcolando l'ampiezza
differenziale come
$$f_k(\vartheta) = -\frac{m\sqrt{2\pi}}{\hbar}\tilde V(\mathbf{q}),\phantom{x}\mathbf{q}=
\mathbf{k'}-\mathbf{k}$$
in coordinate sferiche, la trasformata di $V$ vale
\begin{gather*}
\tilde V(\mathbf{q}) = \left(\frac{1}{2\pi}\right)^{\frac{3}{2}}\int d^3\mathbf{x}
e^{-i\mathbf{q}\cdot\mathbf{x}}V(\mathbf{x}) = \frac{A}{\left(2\pi\right)^{\frac{3}{2}}}
\int_{0}^{2\pi}d\varphi\int_{0}^{\infty}dr\int_{-1}^{1}d\cos\vartheta r^2 e^{-iqr\cos\vartheta}
\frac{e^{-\mu r}}{r}=\\
=\frac{A}{\sqrt{2\pi}}\int_{0}^{\infty}dr re^{-\mu r}\frac{e^{-iqr}-e^{iqr}}{-iqr} =
\frac{2A}{\sqrt{2\pi}} \frac{1}{q^2+\mu^2}
\end{gather*}
dove nell'ultimo passaggio si è integrato utilizzando il lemma di Jordan e il teorema dei residui.\\
Si ha dunque
$$f_k(\vartheta) = -\frac{2mA}{\hbar^2}\frac{1}{4k^2\sin^2\left(\frac{\vartheta}{2}\right)+\mu^2}
\rightarrow \boxed{\dv[]{\sigma}{\Omega}(\vartheta) = |f_k(\vartheta)|^2 = \frac{4m^2A^2}{\hbar^4}
\frac{1}{\left(4k^2\sin^2\left(\frac{\vartheta}{2}\right)+\mu^2\right)^2}}$$
\textbf{Oss.} Per $\mu\to0$ e $A=zZe^2$ il potenziale si riduce a quello Coulombiano
e la sezione d'urto si riduce a quella di Rutherford: questo è dovuto al fatto che $A$ è la
\textit{costante di accopiamento} dell'interazione e $\mu$ dipende dalla \textit{massa del mediatore},
che, nel caso dell'interazione elettromagnetica, è un fotone con massa $m_\gamma=0$.
\section{Diffusione in onde parziali}
Si considera ora un'approssimazione diversa da quella di Born, applicabile solo ai potenziali centrali
e detta \textit{decomposizione in onde parziali}: l'idea di fondo è sfruttare il SCOC
$\left\{L^2,L_z,H\right\}$ al posto delle tre componenti del momento $\left\{p_x,p_y,p_z\right\}$ per
descrivere il sistema.

Si noti che ogni funzione d'onda di particella libera con $k$ fissato può essere scritta
nella forma
$$\psi_k(r,\vartheta,\varphi)=\sum_{l,m}X_{klm}\varphi^0_{klm}(r,\vartheta,\varphi)$$
dove $\varphi^0_{klm}(r,\vartheta,\varphi)$ sono le autofunzioni simultanee del SCOC 
$\left\{L^2,L_z,H_0\right\}$, dette
\textit{onde sferiche}.

Qualitativamente, si cercano problemi in cui questa serie non abbia troppi termini significativi:
è noto che il parametro di impatto $b$ (distanza tra la "traiettoria" della particella libera e il
centro diffusore) classicamente determini quando la particella interagisce ($b<A$); considerando
$|L|=|pb|$, la particella interagisce solo per $|L|<L_\text{max} = pA$: passando agli autovalori di $L^2$,
la particella interagisce solo se
$$\hbar\sqrt{l(l+1)}\approx\hbar l<pA \leftrightarrow l<kA$$
l'approssimazione quindi, a spanne, converrà solo quando $k$ è piccolo, perché ci saranno meno termini
nella serie.
\subsection{Particella libera}
L'equazione di Schrödinger della particella libera con energia $E=\frac{\hbar^2k^2}{2m}$ è
$$H\psi = E\psi\rightarrow \laplacian\psi - k^2\psi = 0$$
in coordinate sferiche si ha che $\laplacian = \frac{1}{r^2}\pdv[]{}{r}\left(r^2\pdv[]{}{r}\right)-
\frac{L^2}{\hbar^2r^2}$ e la soluzione è fattorizzabile come
$$\psi(r,\vartheta,\varphi) = \varphi^0_{kl}(r)Y^{lm}(\vartheta,\varphi):=
\varphi^0_{klm}(r,\vartheta,\varphi)$$
dove $Y^{lm}$ sono le note armoniche sferiche, autofunzioni simultanee di $L^2$ e $L_z$, e la funzione
radiale $\varphi^0_{kl}(r)$ è tale che la sua funzione ridotta $\chi^0_{kl}:=\frac{\varphi^0_{kl}}{r}$
soddisfi \textit{l'equazione radiale ridotta}
$$(\chi^{0}_{kl})'' + \left(k^2-\frac{l(l+1)}{r^2}\right)\chi^0_{kl} = 0$$
L'andamento di queste funzioni per $r\to\infty$ è
\begin{align*}
    (\chi^0_{kl})'' + k^2\chi^0_{kl}=0&\rightarrow\chi_{\pm kl}^0(r)\sim e^{\pm ikr}\\
				      &\rightarrow\varphi^0_{\pm kl}(r)\sim \frac{e^{\pm ikr}}{r}
\end{align*}
Per risolvere invece l'equazione $\forall r$ conviene utilizzare l'equazione per $\varphi_{kl}^0$ nella
variabile adimensionale $\rho=kr$, riscritta come segue
$$\rho^2(\varphi^0_{kl})''(\rho)+2\rho(\varphi^0_{kl})'(\rho)+(\rho^2 - l(l+1))\varphi^0_{kl}(\rho)=0$$
questa è \textbf{l'equazione di Bessel sferica}.
\subsubsection{Equazione di Bessel}
Si tratta di una equazione differenziale del tipo
$$z^2u'' + zu' +(z^2-l^2)u=0$$
che presenta una singolarità fuchsiana in $z=0$ con indici $\pm l$ e una singolarità irregolare in
$z=\infty$. È una ODE del secondo ordine, quindi presenta due soluzioni linearmente indipendenti:
per $l\not\in\mathbb{Z}$ la soluzione generale è
$$u(z) = AJ^l+BJ^{-l},\text{ con } J^l(z) = \sum_{n=1}^{\infty}\frac{(-1)^n}{\Gamma(n+l+1)n!}
\left(\frac{z}{2}\right)^{2n+l}
$$
dove $J^l(z)$ è detta \textit{funzione di Bessel di I specie}. Si può mostrare che per $l\in\mathbb{Z}$
le due soluzioni non sono indipendenti, cioè $J^l\propto J^{-l}$: in questo caso si usa un'altra
soluzione dell'ODE, la \textit{funzione di Von Neumann} $Y_l(z)$.

L'equazione di Bessel sferica ha una forma leggermente diversa
$$z^2u'' + 2zu' +(z^2-l(l+1))u=0$$
ma le sue soluzioni sono legate a $J^l$ e $Y_l$ nel modo seguente:
\begin{align*}
&\text{Funzione di Bessel sferica} \rightarrow j_l(z) =\sqrt{\frac{\pi}{2z}}J^{l+\frac{1}{2}}(z)\\
&\text{Funzione di Von Neumann sferica} \rightarrow \eta_l(z) =\sqrt{\frac{\pi}{2z}}Y_{l+\frac{1}{2}}(z)=
(-1)^{l+1}\sqrt{\frac{\pi}{2z}}J^{-l-\frac12}(z)
\end{align*}
Partendo dalla loro definizione in funzione della $J^l$, si possono ricavare gli sviluppi asintotici
per $z\to0,\infty$:
\begin{align*}
&j_l(z)\sim
\begin{dcases}
    \frac{1}{(2l+1)!!}z^l,&z\to0\\
    \frac{1}{z}\sin\left(z-l\frac{\pi}{2}
\right),&z\to\infty
\end{dcases}\\\phantom{x}\\
&\eta_l(z)\sim
\begin{dcases}
    -\frac{(2l-1)!!}{z^{l+1}},&z\to0\\
    -\frac{1}{z}\cos\left(z-\frac{l\pi}{2}\right),&z\to\infty
\end{dcases}
\end{align*}
Le funzioni sferiche, per $l\in\mathbb{N}$, sono esprimibili come derivate di funzioni elementari nel modo
seguente:
\begin{align*}
\begin{dcases}
j_0(z) = \frac{\sin z}{z}\\
\eta_0(z)=-\frac{\cos z}{z}
\end{dcases} \longrightarrow
\begin{dcases}
j_l(z) = (-1)^lz^l\left(\frac{1}{z}\dv[]{}{z}\right)^lj_0(z)\\
\eta_l(z)= (-1)^{l+1}z^l \left(\frac{1}{z}\dv[]{}{z}\right)^l\eta_0(z)
\end{dcases}
\end{align*}

Tornando al problema della particella libera, la soluzione dell'equazione è dunque
$$\varphi^0_{kl}(\rho) = Aj_l(\rho)+B\eta_l(\rho)$$
La $\eta_l$ diverge in 0 ($l+1>0$), quindi $B=0$; in $r\to\infty$ la soluzione diventa
$$\varphi^0_{kl}(kr) \underset{r\to\infty}\sim\frac{1}{kr}\frac{1}{2i}
\left(e^{-i\frac{l\pi}{2}}e^{ikr}-e^{i\frac{l\pi}{2}}e^{-ikr}\right) =
e^{-i\frac{l\pi}{2}}\frac{1}{2ikr}\left(e^{ikr}-e^{il\pi}e^{-ikr}\right)$$
che è coerente con l'andamento asintotico dell'equazione di partenza
$\varphi^0_{kl}\sim\frac{e^{\pm ikr}}{r}$.

\textbf{Oss.} Questo andamento asintotico evidenzia che a $r\to\infty$ la funzione d'onda sia
composta dalla somma di un'onda sferica entrante e una uscente l'origine, sfasate di $l\pi$: l'onda
sferica entrante viene riflessa dal centro diffusore, mantenendo il numero d'onda $k$ ma sfasandola.

La soluzione del problema 3D libero è dunque
$$\varphi^0_{klm}(r,\vartheta,\varphi) = \sqrt{\frac{2k^2}{\pi}}j_l(kr)Y^{lm}(\vartheta,\varphi)$$
\textbf{Oss.} Si consideri la densità di probablilità radiale associata a $\varphi^0_{kl}$:
$$p(r)\propto\rho^2|j_l(\rho)|^2\underset{\rho\to0}\sim\rho^{2l+2}$$
allora si può affermare che, approssimativamente
$$p(r)\approx 0\phantom{x} \forall\rho=kr\leq\sqrt{l(l+1)}$$
cioè che l'onda sferica si "ferma" tanto più lontano dal centro diffusore quanto $l$ è grande:
dunque, per avere contributi da pochi termini, occorre che la maggioranza dei valori di $l$ sia
tale che l'onda sferica si fermi prima della zona di influenza del potenziale.
